In this section we present \replay{}, an algorithm for computing metagradients
of large-scale iterative ML algorithms. We first detail the setting, 
then discuss existing approaches to computing metagradients, 
and conclude by describing \replay{}.

\subsection{What is a metagradient?}
\label{subsec:meta_gradient}
Training a machine learning model is a two-step process. 
First, we decide on a {\em training setup}---we must pick,
for example, a neural network architecture, 
a training dataset, and an optimizer for training.
Second, we apply the algorithm defined by this training setup to 
train a model.

Our overall goal in this paper is to optimize model behavior as a 
function of the training setup 
(or, as we call it, the {\em metaparameters}) using
gradient-based methods. To this end, we define the following notation:
\begin{itemize}
	\item Let $\mathbf{z} \in \mathbb{R}^n$ be a vector of continuous 
    metaparameters representing the aspects of the training setup we aim to optimize. 
    For example, if we only want to adjust the learning rate and weight decay of
    SGD then $n = 2$. We handle discrete metaparameters (e.g., choice of training data)
    by finding a continuous relaxation (e.g., importance weights).
	\item Let $\mathcal{A}$ be an {\em algorithm} mapping $\mathbf{z}$ to a trained
    machine learning model; we assume all other
	aspects of the training setup outside $\mathbf{z}$ are fixed and thus part of the
    algorithm $\mathcal{A}$.  
    \item Finally, let $\phi$ be an {\em output function} mapping a model
    $\theta$ to a vector $\phi(\theta) \in \mathbb{R}$.
    For example, $\phi(\theta)$ might represent the validation loss of the model $\theta$.
    We require that $\phi$ be differentiable with respect to $\theta$, but otherwise 
    make no assumptions on $\phi$.
\end{itemize}
With this notation in place, we define the 
{\em training function} $f := \phi \circ \mathcal{A}$
mapping the training setup $\mathbf{z}$ 
{\em directly} to the output function $\phi$ evaluated on the corresponding model. 

Finally, the {\em metagradient} is the gradient of the training function with 
respect to the metaparameters, $\nabla_\mathbf{z} f(\mathbf{z})$. 
Intuitively, the metagradient defines the ``direction of steepest ascent'' in
metaparameter space.


\paragraph{Our focus: iterative algorithms.}  
To efficiently compute the metagradient, 
we restrict our focus to cases where the algorithm $\mathcal{A}$ is 
{\em iterative}, i.e., when it can be written in the form
\begin{align}
  \label{eq:iterative_algo}
  \underbrace{\mathcal{A}(z) \coloneqq \st{T}}_{\mathclap{\text{model state after $T$ steps}}}, 
  \quad \text{where} \quad \underbrace{\st{t+1} := h_t(\st{t}, \mathbf{z})}_{\mathclap{\text{optimizer step $t$}}}.
\end{align}
Here, $\st{t}$ is the optimizer state at step $t$
(with $\st{0}$ being the initial state)
and $h_t$ is the {\em update} mapping from state $t$ to state $t+1$.
The form of~\eqref{eq:iterative_algo} captures most large-scale training algorithms. 
For example, if the setup $\mathbf{z} \in \mathbb{R}^T$ is a \textit{per-step} learning rate, 
and the algorithm $\mathcal{A}$ is 
full batch gradient descent, then each update $h_t$ is
$$
  h_t(\st{t}, \mathbf{z}) \coloneqq \st{t} - z_t \nabla \ell(\st{t}),
$$
where $z_t$ is the learning rate at step $t$, 
$\ell$ is the training loss, and 
the state $\st{t}$ comprises the parameters at step~$t$.
For more complex algorithms like
Adam~\citep{kingma2015adam}, the state $\st{t}$ includes
terms like gradient moments.

\subsection{Warmup: Metagradients via autodifferentiation}
A key primitive we leverage to calculate metagradients is
{\em automatic differentiation} (AD)---a standard tool for taking 
gradients through computer-defined functions.
AD takes gradients 
by decomposing functions into elementary operations 
with known derivatives, then combining these derivatives using the chain rule.
Concretely, AD operates in two passes: a ``forward pass,'' 
which executes the function of interest and stores intermediate products 
for {each} elementary operation; and a ``backward pass,''
which calculates the gradient by propagating chains of
partial derivatives using these stored products. For the purposes of this paper, 
we will view AD as a black box that calculates the gradient of a many-to-one function 
(i.e., any $f: \mathbb{R}^d \to \mathbb{R}$)
at a given point using only a small constant factor more time than calculating the
function itself (along with the space cost of storing the 
necessary forward-pass products).

What does this have to do with metagradients? 
Well, seeing as how training itself is a computer-defined function, 
AD is a natural tool for calculating the metagradient. The main challenge, 
as we discuss in the sequel, is that AD-based approaches to calculating 
the metagradient tend to be too resource-intensive for the large-scale machine
learning algorithms we consider.
In the remainder of this section we build up background before finally
describing \replay{}, our algorithm for scalably computing (exact)
metagradients.

\paragraph{Approach \#1: Direct AD.} The direct approach to calculating
metagradients exploits the fact that nearly any learning algorithm
is itself a sequence
of differentiable computer-defined operations---meaning the training function $f$
is \emph{also differentiable}.

However, operationalizing this observation to compute metagradients
turns out to be challenging.
The reason is that AD stores intermediate
products for {\em each} operation.
The amount of data stored thus scales with the number of operations 
in the function of interest. In the case of our training function $f$,
this number encompasses {\em all} the operations used to train a machine 
learning model.
As a result, even in a toy scenario like MNIST training, computing metagradients
with na\"ive AD would require storing terabytes of data. 

\paragraph{Approach \#2: Exploiting structure with step-wise AD.} 
A more efficient method for calculating the metagradient,
{\em step-wise AD},
leverages the structure of iterative learning algorithms~\citep{werbos1990backpropagation,maclaurin2015gradient,franceschi2017forward}.
Recall from \eqref{eq:iterative_algo} that such algorithms take the form
\[
	{\mathcal{A}(\mathbf{z}) \coloneqq \st{T}}, \quad \text{where} \quad {\st{t+1} := h_t(\st{t}, \mathbf{z})}.
\]
Algebraic manipulation 
(in particular, using the chain rule, 
the law of the total derivative,
and the identity $\st{t} = h_{t-1}(\st{t-1}, \mathbf{z})$)
allows us to write the metagradient over an iterative algorithm as
\begin{equation}
	\pfrac{f(\mathbf{z})}{\mathbf{z}} =
	\pfrac{\phi(\mathcal{A}(\mathbf{z}))}{\mathbf{z}} 
	= \mathlarger{\sum}_{t=1}^{T}\ 
  \underbrace{\overbrace{\pfrac{\phi(\st{T})}{\st{t}}}^{A_t} \cdot 
  \pfrac{h_{t-1}(\st{t-1}, \mathbf{z})}{\mathbf{z}}}_{B_t} \label{eq:total_deriv}, %
\end{equation}
where we have introduced the notation $A_t$ and $B_t$ for notational convenience.
Step-wise AD computes the metagradient by calculating each term in the sum of
\eqref{eq:total_deriv} one at a time. 
For each term, the main challenge lies in computing $A_t$, since given $A_t$ 
we can straightforwardly compute $B_t$ (the entire term)
by differentiating through a single model update, i.e.,
\begin{align*}
  B_t := A_t \cdot \pfrac{h_{t-1}(\st{t-1}, \mathbf{z})}{\mathbf{z}}
  =  \pfrac{(A_t \cdot h_{t-1}(\st{t-1}, \mathbf{z}))}{\mathbf{z}},
\end{align*}
which is just a single call to our assumed ``AD oracle'' on the function 
$\mathbf{z} \mapsto A_t \cdot h_{t-1}(\st{t-1}, \mathbf{z})$.
Computing the $A_t$ terms is less straightforward as we need to relate $s_t$ and $s_T$;
to do so, we exploit the recurrence
\begin{equation}
  A_t := \pfrac{\phi(\st{T})}{\st{t}} 
  = \pfrac{\phi(\st{T})}{\st{t+1}} \cdot \pfrac{{h_t(\st{t}, \mathbf{z})}}{\st{t}}
  = \pfrac{(A_{t+1} \cdot h_t(\st{t}, \mathbf{z}))}{\st{t}},\label{eq:grad_recurrence}
\end{equation}
making $A_t$ straightforward to compute (again, a single ``AD oracle'' call) given $A_{t+1}$. 
Step-wise AD exploits this fact to successively
calculate the gradient with respect to each state, from state $T$ down to state $0$.

Bringing these ingredients together, the algorithm executes as follows. As a
preprocessing step, it trains the model and stores all intermediate states
$\st{0},\ldots,\st{T}$.
Then, the algorithm calculates and sums the terms in
\eqref{eq:total_deriv}. It first computes $A_T := \nicefrac{\partial
\phi(\st{T})}{\st{T}}$, the gradient of the output function $\phi$ with respect
to the final state. Then, the algorithm steps through $\st{T-1},\ldots,\st{0}$ in
reverse order, calculating (a) the gradient with respect to each
state $A_t$ (via \eqref{eq:grad_recurrence}) and (b) the gradient with respect to
$\mathbf{z}$ at that step $B_t$ (via \eqref{eq:total_deriv}, using the
previously calculated gradient with respect to that state). AD calculates
both quantities---each requires differentiating over only one train step.
Finally, the algorithm returns the final metagradient as the sum of the terms.

Despite improving storage overhead compared to ``direct AD'', step-wise
AD is still too space-intensive at scale. After all,
this algorithm saves {\em every} optimizer state. 

\begin{figure}
	\onecolumn
    \centering
    \begin{forest}
    for tree={
      fit=band,
      minimum size=2.2em,
      s sep=0.46em,
    rectangle,
    fit=band,
    draw,
    calign primary child=1,
    calign=child edge
    },
    [{$0$}, simple cylinder, fill=gray!20, name=el1 %
      [{$0$}, simple cylinder, fill=gray!20
        [$\cdots$,circle,draw=none
          [{$0$}
            [{$0$}]
            [{$1$}]]
          [$\cdots$,circle,draw=none
            [$\cdots$,circle,draw=none]]]
        [$\cdots$,circle,draw=none
          [$\frac{n}{4}$
            [$\frac{n}{4}$] [$\frac{n}{4}+ 1$]]
          [$\cdots$,circle,draw=none
            [$\cdots$,circle,draw=none]]]]
      [$\frac{n}{2}$, simple cylinder, fill=gray!20, name=el2
        [$\cdots$,circle,draw=none,
          [$\frac{n}{2}$
            [$\frac{n}{2}$]
            [$\frac{n}{2}+ 1$]]
          [$\cdots$,circle,draw=none
            [$\cdots$,circle,draw=none]]]
        [$\cdots$,circle,draw=none, name=el3
          [$\frac{3n}{4}$, simple cylinder, fill=gray!20,name=el8
            [$\frac{3n}{4}$, simple cylinder, fill=gray!20]
            [$\frac{3n}{4} \footnotesize{+1}$, simple cylinder, minimum height=1em, minimum size=2em, text height=.75em, text depth=3.4pt, fill=gray!20,name=el9]]
          [$\cdots$,circle,draw=none,name=el35
            [$\cdots$,circle,draw=none,name=el5]
            [{$n\shortminus{}1$},name=el4]]]]]
    \begin{scope}[>=Triangle]
        \path [draw, <->] (current bounding box.south west) +(-20pt,0) coordinate (c) -- (c |- current bounding box.north) node [pos=.5, fill=white] {$\log_2(n)$};
    \end{scope}
    \draw[->,thick,color=blue] (el1) to[out=40,in=110] node [pos=0.40, fill=white, yshift=-4pt] {Traversal order} (el2); %
    \draw[->,thick,color=blue] (el2) to[out=40,in=110] (el3);
    \draw[->,thick,color=blue] (el3) to[out=40,in=110] (el35);
    \draw[->,thick,color=blue] (el35) to[out=east,in=north] (el4);
    \draw[->,thick,color=blue] (el4) to[out=south,in=south] (el5);
    \draw[->,thick,color=blue] (el5) to[out=north west,in=south west] (el35);
    \draw[->,thick,color=blue] (el35) to[out=west,in=south east] (el3);
    \draw[->,thick,color=blue] (el3) to[out=292,in=70] (el8);
    \draw[->,thick,color=blue] (el8) to[out=east,in=north] (el9);
    \begin{scope}[>=Triangle]
    \path [draw, <->] (current bounding box.south west) +(43.75pt,-1pt) coordinate (c) -- (current bounding box.south east) node [pos=.5, fill=white] {$n$ states}; %
    \end{scope}
    \end{forest}

    \caption{The lazy $k$-ary tree structure for
    traversing optimizer states in reverse order, with $k=2$. Recall that $n$ is
    the number of states (parameterized such that $n=T+1$). Each node represents
    the correspondingly numbered state. We give an example of the traversal
    using the \color{blue}{blue arrows }\color{black} in the figure, which
    denote the traversal path up to state $\smash{s_{\frac{3n}{4} + 1}}$. The gray
    cylinders \hspace{1pt}\tikzinlinecylinder{}\hspace{1pt} indicate the states
    that are stored when the traversal is at state $\smash{s_{\frac{3n}{4} + 1}}$; the
    other states are not stored at this point in the traversal. Traversing this
    structure requires storing $\mathcal{O}(\log(n))$ state and computing
    $\smash{\mathcal{O}(n\log(n))}$ optimizer steps---compared to $n$ for simply
    training. }
    \label{fig:efficient_data_structure}
\end{figure}

\subsection{\replay{}} 
\replay{} is our algorithm for efficiently and exactly computing metagradients.
It uses $\mathcal{O}(k\log_k(T))$ space and requires running the learning algorithm $\mathcal{A}$
a total of $1 + \log_k(T)$ times, 
with $k$ a user-chosen constant. The main idea is to make the space-intensive subroutine of
step-wise AD---a reverse-order traversal of the optimizer states at each
step---much more efficient. After all, step-wise AD stores \textit{all} the
states to reverse traverse them. \replay{} modifies step-wise AD to traverse
states in less space by exploiting a simple observation: when training is
deterministic, one can \textit{reinstantiate} an optimizer state $\st{t}$ by
``replaying'' training from a fixed point $t' < t$---at the compute cost of $t -
t'$ training steps. For example, one simple scheme saves every other state, then
``replays'' the remaining states when (reverse) traversing; this routine
stores $T/2$ states but computes an extra $T/2$ model updates compared
to storing \textit{all} the states.

\replay{} performs a reverse-order traversal the optimizer states while
balancing  the compute cost of ``replaying'' training with the storage cost of
saving states. We use a combination of deterministic training 
(fixing data ordering, data augmentation, and any other randomness in 
the training process) and an efficient data structure (similar to a segment
tree; see Figure \ref{fig:efficient_data_structure}) 
to reverse-order traverse the optimizer states with
$\mathcal{O}(k\log_k(T))$ space and an additional $T\log_k(T)$
model steps. 

Specifically, \replay{} recursively saves and replays training states. The
algorithm splits the training trajectory into $k$ segments, performs the full
training routine while saving only the start of each segment, then recurses
into each segment (in reverse) to retrieve the states in reverse-order. The
recursion depth bottoms out at $\log_k(T)$, at which point the algorithm has
$k$ consecutive optimizer states in memory; the algorithm then backpropagates
along this segment, before deleting all these states from memory and then
reinstantiating the next $k$-length segment of optimizer states. 
We provide additional details on the algorithm in Appendix~\ref{app:replay}.
\replay{} unlocks computing large-scale metagradients by requiring only
logarithmic storage and additional compute time. 



\begin{remark}[Connection to rematerialization]
In a broad sense, both \replay{} and step-wise AD above can be viewed as
special cases of a classical approach in AD (and computing broadly) known as
\textit{rematerialization}~\citep{chaitin1981register,briggs1992rematerialization,zweig2000exact,griewank2008evaluating,chen2016training}.
To our knowledge, however, \replay{} is the first application of this 
particular rematerialization technique to the problem of computing metagradients through 
model training.
\end{remark}

\begin{remark}[Reversible learning]
  An alternative approach to calculating metagradients that does not save any 
  state is \underline{\smash{reversible}} learning~\citep{maclaurin2015gradient},
  for which one can ``invert'' previous training states from future ones.
  We focus here on general (non-reversible) learning algorithms for two reasons: first, 
  even simple algorithms such as SGD without momentum are non-reversible; second, 
  reversibility in practice introduces numerical precision issues.
\end{remark}
